\documentclass[12pt]{article}
\usepackage[utf8]{inputenc}
\usepackage[T1]{fontenc}
\usepackage[margin=1in]{geometry}
\usepackage{setspace}
\usepackage{xcolor}
\usepackage[colorlinks=true, urlcolor=blue!60!black, linkcolor=blue!60!black]{hyperref}
\setlength{\parindent}{0em}
\setlength{\parskip}{0.45em}
\renewcommand{\baselinestretch}{1.0}
\begin{document}
\begin{center}
{\Large \textit{Seminole Tribe of Florida v.\ Florida}, 517 U.S.\ 44 (1996)\par}
\vspace{0.3em}
{\normalsize Case Brief --- One-Page Edition\par}
\end{center}

\textbf{Memory Jogger.} No Article I abrogation of state sovereign immunity; \textit{Union Gas} overruled (\href{https://www.law.cornell.edu/supct/html/94-12.ZO.html}{Cornell LII}).

\textbf{Facts.} The Indian Gaming Regulatory Act required States to negotiate gaming compacts in good faith and let tribes sue them in federal court. The Seminole Tribe sued Florida and Governor Chiles for refusing to negotiate.

\textbf{Procedural History.} The district court upheld abrogation; the Eleventh Circuit ordered dismissal.

\textbf{Votes.} Affirmed, 5--4: Rehnquist for the Court; Stevens and Souter (with Ginsburg and Breyer) dissenting.

\textbf{Holding.} Article I powers, including the Indian Commerce Clause, cannot abrogate state immunity; Section~5 of the Fourteenth Amendment can. A detailed remedial scheme displaces \textit{Ex parte Young}.

\textbf{Reasoning.} Under \textit{Hans}, immunity is a constitutional ``presupposition''; \textit{Fitzpatrick} rested on the later Fourteenth Amendment, and \textit{Union Gas} drew no principled line. Stevens called the ruling an ``affront to a coequal branch''; Souter said the Amendment repealed only diversity jurisdiction.

\textbf{Examples.} \emph{Same-Side:} copyright damages suits against States. \emph{Opposite-Side:} Section~5 family-leave suits. \emph{Fence-Sitter:} Bankruptcy Clause suits.

\textbf{Critique.} Textualists embraced an atextual immunity, and federal rights lost effective remedies.

\textbf{Key Quotation.} ``Even when the Constitution vests in Congress complete lawmaking authority over a particular area, the Eleventh Amendment prevents congressional authorization of suits by private parties against unconsenting States. The Eleventh Amendment restricts the judicial power under Article III, and Article I cannot be used to circumvent the constitutional limitations placed upon federal jurisdiction.'' --- Chief Justice Rehnquist.

\end{document}
